\begin{figure}[tbp]
\centering
\begin{tikzpicture}[x=1cm,y=1cm,
  witness/.style={font=\small,inner sep=2pt},
  inclusion/.style={densely dashed},
  panel/.style={draw=black!30,rounded corners=3pt,inner sep=7pt}]
\node[witness] (h) at (0,0) {\(H\)};
\node[witness] (t1) at (-2.4,1.05) {\(T_1\)};
\node[witness] (t2) at (2.4,1.05) {\(T_2\)};
\node[witness] (k11) at (-3.7,2.2) {\(K_{X_1}\)};
\node[witness] (k12) at (-2.4,2.2) {\(K_{X_2}\)};
\node[witness] (k1q) at (-1.1,2.2) {\(K_{X_q}\)};
\node[witness] (k21) at (1.1,2.2) {\(K_{Y_1}\)};
\node[witness] (k22) at (2.4,2.2) {\(K_{Y_2}\)};
\node[witness] (k2q) at (3.7,2.2) {\(K_{Y_q}\)};
\node[font=\scriptsize] at (-1.75,2.2) {\(\cdots\)};
\node[font=\scriptsize] at (3.05,2.2) {\(\cdots\)};
\foreach \target in {k11,k12,k1q}{\draw (t1)--(\target);}
\foreach \target in {k21,k22,k2q}{\draw (t2)--(\target);}
\draw[inclusion] (h)--(t1);
\draw[inclusion] (h)--(t2);
\node[panel,fit=(t1)(k11)(k12)(k1q)] (p1) {};
\node[panel,fit=(t2)(k21)(k22)(k2q)] (p2) {};
\node[font=\small,above=3pt of p1] {First private chart};
\node[font=\small,above=3pt of p2] {Second private chart};
\node[font=\small,below=7pt of h] {\(K_{X_i}\cap K_{Y_j}=H\quad\text{for every }i,j\)};
\end{tikzpicture}
\caption{The subgroup relations supplied by \Cref{c:failed-pattern}.
In the first chart, the witnesses on one parameter line satisfy
\(\Sigma_1=\langle X_i,X_j\rangle\) for every \(i\ne j\),
\(\operatorname{rk}\Sigma_1=\operatorname{rk}X_i+1\), and
\(T_1=K_{\Sigma_1}=K_{X_i}\cap K_{X_j}\); each \(K_{X_i}\) covers \(T_1\).
The second chart has the analogous relations with \(Y_j,\Sigma_2,T_2\).
Solid segments denote these covers, and dashed segments denote inclusions,
with intervening labels omitted. The diagram is schematic in rank and
shows only the stated relations.}
\label{c:pattern-figure}
\end{figure}
